\section{Symmetry in quantum mechanics}

\referto{\tong~7.1,7.2.1,7.5\\
\sakurai~4.1, Goldstein et al ``Classical Mechanics'' 3.9 for Laplace-Runge-Lentz vector\\ 
\sakurai~4.4, Weinberg ``Quantum theory of fields'' Vol 1 Chapter 2.2 and Appendix 2.A for Wigner's theorem on symmetry operations}

\subsection{Symmetry in Classical Mechanics}

A continuous symmetry keeps the Lagrangian invariant:
\begin{equation}
 q_i\longmapsto q_i'=q_i+\epsilon\,\delta q_i,
 \qquad
 \mathcal L(q_i',\dot q_i')=\mathcal L(q_i,\dot q_i).
\end{equation}
Using
\begin{equation}
 \frac{\partial\mathcal L}{\partial q_i}
 =\frac{\dif}{\dif t}
 \frac{\partial\mathcal L}{\partial\dot q_i},
\end{equation}
one obtains
\begin{equation}
 \sum_i\left(
 \delta\dot q_i\frac{\partial\mathcal L}{\partial\dot q_i}
 +\delta q_i\frac{\partial\mathcal L}{\partial q_i}
 \right)
 =\frac{\dif Q}{\dif t}=0,
\end{equation}
where
\begin{equation}
 \boxed{Q=\sum_i\delta q_i
 \frac{\partial\mathcal L}{\partial\dot q_i}}
\end{equation}
is the conserved Noether charge.

\subsubsection{Translations}
For
\begin{equation}
 \mathcal L=\frac12m\dot{\bm r}^{\,2}+\frac{k}{r},
\end{equation}
a translation has
\begin{equation}
 \delta\bm r=\bm a,
 \qquad \bm a=\text{constant}.
\end{equation}
The Noether charge is
\begin{equation}
 \sum_j\delta r_j\frac{\partial\mathcal L}{\partial\dot r_j}
 =\bm a\cdot m\dot{\bm r},
\end{equation}
so the momentum
\begin{equation}
 \boxed{\bm p=m\dot{\bm r}}
\end{equation}
is conserved.

\subsubsection{Rotations}
For $R\in O(3)$,
\begin{equation}
 \bm r'=R\bm r,
 \qquad
 \mathcal L(\bm r',\dot{\bm r}')
 =\mathcal L(\bm r,\dot{\bm r}).
\end{equation}
An infinitesimal rotation in the $y$-$z$ plane is
\begin{equation}
 R(\theta)=
 \begin{pmatrix}
 1&0&0\\
 0&\cos\theta&-\sin\theta\\
 0&\sin\theta&\cos\theta
 \end{pmatrix}
 \simeq
 \begin{pmatrix}
 1&0&0\\
 0&1&-\theta\\
 0&\theta&1
 \end{pmatrix}.
\end{equation}
Thus
\begin{equation}
 \delta y=-\theta z,
 \qquad
 \delta z=\theta y,
\end{equation}
and
\begin{equation}
 \boxed{L_x=-z\frac{\partial\mathcal L}{\partial\dot y}
 +y\frac{\partial\mathcal L}{\partial\dot z}
 =yp_z-zp_y.}
\end{equation}

\subsubsection{Runge--Lenz vector}
For the Coulomb/Kepler problem,
\begin{equation}
 \boxed{\bm A=\bm p\times\bm L-mk\hat{\bm r}.}
\end{equation}
The infinitesimal transformation generated by $A_j$ is
\begin{equation}
 \delta_j r_i
 =\frac{\epsilon}{2}
 \left(2p_jr_i-r_jp_i-\delta_{ij}\bm r\cdot\bm p\right).
\end{equation}
The Lagrangian variation is
\begin{equation}
 \delta_j\mathcal L
 =\epsilon\left(
 m\dot r_i\delta_j\dot r_i-\frac{k}{r^2}\delta_jr
 \right)
 =\epsilon\frac{\dif A_j}{\dif t}.
\end{equation}

\subsection{Implementing symmetry in Quantum Mechanics}

\subsubsection{Wigner's theorem}

For observable 
\begin{equation}
 \hat O=\hat O^\dagger=\sum_q q\proj{q},
\end{equation}
the probability of measurement outcome $q$ in state $\dket\psi$ is given by
\begin{equation}
 \text{Pr}(q)=|\expval{q}{\psi}|^2.
\end{equation}
It's very natural to assume that the measurement probability should remain invariant under a symmetry transformation. Wigner proved the following theorem on how such a symmetry transformation is implemented on states in the Hilbert space $\mathbb{V}$:

\begin{thm}
  \textbf{(Wigner's theorem on symmetry operations)} Consider a bijective symmetry transformation that preserves transition probabilities $\expval{\psi|\phi}$ for any two states $\dket\psi,\dket\phi\in\mathbb{V}$. It is either implemented by a unitary linear operator $\hat U$, or by an anti-unitary anti-linear operator $\hat U\cdot\mathcal{K}$ where $\mathcal{K}$ stands for complex conjugation.   
\end{thm}
Below we discuss these two scenarios in detail. 

\subsubsection{Unitary symmetry}
The symmetry transformation is implemented by a unitary operator
\begin{equation}
 \hat U=\sum_q\dket{\widetilde q}\dbra q,
 \qquad
 \hat U^\dagger\hat U=\id,
\end{equation}
which transforms a complete orthonormal basis set $\{\dket q\}$ into a new set of complete orthonormal basis $\{\dket{\widetilde q}\}$. 

The inner product remains invariant under a unitary symmetry transformation:
\begin{equation}
 \dket{\widetilde\psi}=\hat U\dket\psi,
 \qquad
 \expval{\widetilde q|\widetilde\psi}=\expval{q|\psi}.
\end{equation}
The transformation is linear:
\begin{equation}
 \hat U(\alpha\dket\psi+\beta\dket\phi)
 =\alpha\hat U\dket\psi+\beta\hat U\dket\phi.
\end{equation}
States and operators transform as
\begin{equation}
 \dket\psi\longmapsto\dket{\widetilde\psi}=\hat U\dket\psi,
 \qquad
 \hat O\longmapsto\hat{\widetilde O}=\hat U\hat O\hat U^\dagger.
\end{equation}
so that the expectation value of observabl $\hat O$ remains invariant
\begin{align}
    \expval{\hat O}=\expval{\psi|\hat O|\psi}=\expval{\widetilde\psi|\hat{\widetilde O}|\widetilde\psi}
\end{align}

\subsubsection{Antiunitary symmetry}
Let $\ck$ denote complex conjugation. An anti-unitary symmtry is implemented by 
\begin{equation}
 \dket{\widetilde\psi}=\hat U \ck\dket\psi=\hat U\dket{\psi^\ast},
 \qquad
 \expval{\widetilde q|\widetilde\psi}
 =\expval{q|\psi}^*=\expval{\psi|q}.
\end{equation}
The transformation is antilinear:
\begin{equation}
 \hat U \ck(\alpha\dket\psi+\beta\dket\phi)
 =\alpha^*\hat U\ck\dket\psi
 +\beta^*\hat U\ck\dket\phi.
\end{equation}
States and operators transform as
\begin{equation}
 \dket\psi\longmapsto\hat U\dket{\psi^*},
 \qquad
 \hat O\longmapsto\hat U\hat O^*\hat U^\dagger.
\end{equation}
A simplest example of anti-unitary symmetries is a quantum particle without eletromagnetic fields:
\begin{equation}
 \hat T=\ck,
 \qquad
 \hat H=\frac{\hat p^2}{2m}+V(\hat{\textbf{r}})=\hat H^\ast.
\end{equation}

\subsection{Examples of Unitary Symmetries}

\subsubsection{Parity in one dimension}
For a quantum particle in one dimension, the inversion (or parity) symmetry is defined as 
\begin{equation}
 \hat I=\int\dif x\,\dket{-x}\dbra x.
\end{equation}
with 
\begin{equation}
 \hat I=\hat I^\dagger,
 \qquad
 \hat I^2=\id.
\end{equation}
If the Hamiltonian preserves inversion symmetry
\begin{equation}
 \comm{\hat I}{\hat H}=0,
\end{equation}
then $\hat I$ and $\hat H$ can be diagonalized simultaneously:
\begin{equation}
\hat H\dket\psi=E\dket\psi,~~~\hat I\dket\psi=\pm\dket\psi.
\end{equation}
Therefore,
\begin{equation}
 \expval{x|\hat I|\psi}=\expval{-x|\psi}
 =\pm\expval{x\psi},
\end{equation}
so each energy eigenfunction can be labeled by its even or odd parity.

\subsubsection{Continuous translations}
Translation by distance $a$ is implemented by
\begin{equation}
 \boxed{\hat U_{\mathrm{trans}}(a)
 =e^{-\imth\hat pa/\hbar}.}
\end{equation}
If
\begin{equation}
 \comm{\hat U_{\mathrm{trans}}(a)}{\hat H}=0,
\end{equation}
then
\begin{equation}
 \comm{\hat p}{\hat H}=0.
\end{equation}
and hence each energy eigenstate is labeled by its momentum eigenvalue.\\
More generally, a continuous unitary symmetry is implemented by
\begin{equation}
 \hat U(\epsilon)=e^{-\imth\epsilon\hat Q/\hbar},
 \qquad
 \hat Q=\hat Q^\dagger.
\end{equation}
If
\begin{equation}
 \comm{\hat U(\epsilon)}{\hat H}=0,
\end{equation}
then
\begin{equation}
 \boxed{\comm{\hat Q}{\hat H}=0.}
\end{equation}
Thus the generator $\hat Q$ of the continuous symmetry is a conserved observable.

\subsubsection{Phase-space rotations of the harmonic oscillator}
For
\begin{equation}
 \hat H_{\mathrm{SHO}}
 =\frac{\hat p^2}{2m}+\frac12m\omega^2\hat x^2
 =\frac{1}{2m}\left[\hat p^2+(m\omega\hat x)^2\right],
\end{equation}
the transformation
\begin{equation}
 \begin{pmatrix}\hat p\\m\omega\hat x\end{pmatrix}
 \longmapsto
 \begin{pmatrix}
 \cos\theta&-\sin\theta\\
 \sin\theta&\cos\theta
 \end{pmatrix}
 \begin{pmatrix}\hat p\\m\omega\hat x\end{pmatrix}
\end{equation}
leaves the Hamiltonian invariant. One can show this symmetry is implemented by a unitary transformation (see HW10). 


\subsubsection{Rotation in two-dimensional $y$-$z$ plane}
The generator is
\begin{equation}
 \hat L_x=\hat y\hat p_z-\hat z\hat p_y.
\end{equation}
The commutators are
\begin{align}
 \comm{\hat y}{\hat L_x}&=-\imth\hbar\hat z,&
 \comm{\hat z}{\hat L_x}&=\imth\hbar\hat y,\\
 \comm{\hat p_y}{\hat L_x}&=-\imth\hbar\hat p_z,&
 \comm{\hat p_z}{\hat L_x}&=\imth\hbar\hat p_y.
\end{align}
Therefore,
\begin{align}
 e^{-\imth\theta\hat L_x/\hbar}\hat y
 e^{\imth\theta\hat L_x/\hbar}
 &=\cos\theta\,\hat y+\sin\theta\,\hat z,\\
 e^{-\imth\theta\hat L_x/\hbar}\hat z
 e^{\imth\theta\hat L_x/\hbar}
 &=\cos\theta\,\hat z-\sin\theta\,\hat y.
\end{align}
Equivalently,
\begin{equation}
 e^{-\imth\theta\hat L_x/\hbar}
 \begin{pmatrix}\hat x\\\hat y\\\hat z\end{pmatrix}
 e^{\imth\theta\hat L_x/\hbar}
 =
 \begin{pmatrix}
 1&0&0\\
 0&\cos\theta&\sin\theta\\
 0&-\sin\theta&\cos\theta
 \end{pmatrix}
 \begin{pmatrix}\hat x\\\hat y\\\hat z\end{pmatrix}.
\end{equation}
The momentum transforms in the same way.

\textbf{Quantization of angular momentum}
A $2\pi$ rotation gives
\begin{equation}
 e^{-\imth2\pi\hat L_z/\hbar}=\id.
\end{equation}
Hence
\begin{equation}
 \boxed{L_z=n\hbar,
 \qquad n\in\mathbb Z.}
\end{equation}

\subsubsection{Rotations in three dimensions}
The three generators are
\begin{align}
 \hat L_x&=\hat y\hat p_z-\hat z\hat p_y,\\
 \hat L_y&=\hat z\hat p_x-\hat x\hat p_z,\\
 \hat L_z&=\hat x\hat p_y-\hat y\hat p_x.
\end{align}
Symbolically,
\begin{equation}
 \boxed{\hat L_\alpha
 =\epsilon_{\alpha\beta\gamma}
 \hat x_\beta\hat p_\gamma.}
\end{equation}
Since
\begin{equation}
 \comm{\hat L_x}{\hat L_y}\neq0,
\end{equation}
rotationally invariant Hamiltonians with $[\hat L_\alpha,\hat H]=0,~\alpha=x,y,z$ can exhibit degenerate energy levels. For
example,
\begin{equation}
 \hat H=\frac{\hat{\bm p}^{\,2}}{2m}+V(r)
\end{equation}
is rotationally invariant for a central potential $V(r)$.


